\chapter{总结与展望}
\section{工作总结}
伴随计算机硬件性能的持续提升以及图形处理能力的增强，图形模拟在游戏开发、电影制作等众多领域得到了广泛应用。流体同弹塑性固体的高精度交互仿真已经成为研究的热点，物质点法（MPM）在这一场景中虽已取得显著成果，但在算法层面和工程角度仍存在优化空间，如模拟物体碰撞、流固耦合时会产生数值耗散，碰撞边界处会出现空白缝隙，且国内缺乏基于节点系统的物理仿真软件。

针对上述问题，本文深入研究，主要贡献如下：其一，针对传统MPM在统一背景网格上存在数值粘性的局限，通过修改传统MPM流程，引入分解兼容仿射粒子网格（DC - APIC）积分器，在切向方向强制实施自由滑移边界条件，在法向方向实现分离条件，并设计基于相场梯度的分割方法，经多种复杂的流体与弹性固体相互作用仿真，验证了该方法的有效性。其二，针对传统MPM处理物体碰撞产生白色缝隙的问题，提出融合DC - APIC积分器的自适应网格算法，在不增加过多额外开销的同时，有效保留不同尺度的碰撞细节。其三，设计并实现基于节点系统的物理仿真软件，具备程序化编辑和运算系统，支持模型导入、物理仿真、光追渲染等功能，并在软件框架内实现了改进的物质点法算法。通过多个具有挑战性的大规模实验场景，验证了本文方法在生成真实感视觉效果方面的鲁棒性和实用性，同时通过多个验证对比实验进行具体分析，成功提供了能正确模拟流体与固体双向耦合交界面正确分离的仿真方法以及一套节点化物理仿真软件，为相关领域发展提供了有力支撑。 
\section{未来展望}
尽管我们的 DC-APIC 传输方案在处理体积固体的情况下是稳定的，但粒子之间仍然可能会相互交叉，这一问题在基于符号距离函数（SDF）的 DC-APIC 方法中已得到解决，该方法通过在水平集上基于穿透距离施加惩罚力来处理。然而，在基于相场梯度的 DC-APIC 方法中，由于缺乏穿透距离信息，仍然没有可靠的解决方案。在薄片、细丝等情况下使用 DC-APIC 仍然极具挑战性，将我们的工作扩展到支持流体与布料以及毛发之间的相互作用将会很有趣。此外，本文只实现了DC-APIC的显式积分器，没有构建全隐式弹性物质点法强耦合模型。未来工作的另一个有前景的方向是调节不相容的网格节点与粒子之间的切向相互作用，以便更好地考虑摩擦效应。